% Generated draft: requires mathematical and source review.
\documentclass[11pt]{article}
\usepackage[margin=25mm]{geometry}
\usepackage{fontspec}
\setmainfont{Latin Modern Roman}
\usepackage{amsmath,amssymb,mathtools,unicode-math}
\setmathfont{Latin Modern Math}
\usepackage{enumitem,needspace,xurl,hyperref}
\setlength{\parindent}{0pt}
\setlength{\parskip}{5pt}
\setlength{\emergencystretch}{3em}
\newcommand{\N}{\mathbb{N}}
\newcommand{\Z}{\mathbb{Z}}
\newcommand{\Q}{\mathbb{Q}}
\newcommand{\R}{\mathbb{R}}
\newcommand{\C}{\mathbb{C}}
\newcommand{\F}{\mathbb{F}}
\newcommand{\A}{\mathbb{A}}
\newcommand{\PP}{\mathbb{P}}
\newcommand{\E}{\mathbb{E}}
\newcommand{\eps}{\varepsilon}
\newcommand{\dd}{\,\mathrm{d}}
\begin{document}
\section*{UnsolvedMath ID: 30002773}
\textbf{Source problem code:} catalogue entry
\subsection{Definitions and mathematical statement}
Let \(j,k\in\Z_{\geq 1}\). The supplied record calls a simultaneous equipartition of \(j\) masses by \(k\) affine hyperplanes a division into \(2^k\) regions, each containing the fraction \(2^{-k}\) of every mass.

The following makes this condition precise under an explicitly conditional measure interpretation. Suppose that, for each positive integer \(d\), \(\mathcal M_d\) is a specified class of finite nonnegative Borel measures on \(\R^d\), each assigning measure zero to every affine hyperplane. For affine hyperplanes
\[
H_i=\{x\in\R^d:\langle a_i,x\rangle=b_i\},\qquad a_i\in\R^d\setminus\{0\},\quad b_i\in\R,\quad 1\leq i\leq k,
\]
define their sign regions by
\[
R_{\sigma}=\bigcap_{i=1}^k\{x\in\R^d:\sigma_i(\langle a_i,x\rangle-b_i)>0\},\qquad \sigma\in\{-1,1\}^k.
\]
They simultaneously equipartition \(\mu_1,\ldots,\mu_j\in\mathcal M_d\) if
\[
\mu_r(R_{\sigma})=2^{-k}\mu_r(\R^d)
\quad\text{for every }r\in\{1,\ldots,j\}\text{ and every }\sigma\in\{-1,1\}^k.
\]
The boundary condition on the measures ensures that excluding the hyperplanes does not discard any mass. Here the source word ``orthants'' is interpreted as these sign regions; perpendicularity of the hyperplanes is not asserted.

For such a specified family \(\mathcal M=(\mathcal M_d)_{d\geq1}\), define, whenever the minimum exists,
\[
\Delta_{\mathcal M}(j,k)=\min\left\{d\in\Z_{\geq1}:\ \forall(\mu_1,\ldots,\mu_j)\in\mathcal M_d^{\,j},\ \exists H_1,\ldots,H_k\text{ simultaneously equipartitioning them}\right\}.
\]
The record denotes its intended quantity by \(\Delta(j,k)\), but does not specify the admissible class of masses. Thus the measure hypotheses above describe a conditional formulation, not hypotheses verified from the supplied source.

\textbf{Ramos conjecture, as supplied.} The stated known lower bound for \(\Delta(j,k)\) is attained for every pair of positive integers \(j,k\). If \(L(j,k)\) denotes that intended lower bound, this assertion has the schematic form
\[
\forall j,k\in\Z_{\geq1},\qquad \Delta(j,k)=L(j,k).
\]
The supplied record gives no formula for \(L(j,k)\). Consequently, neither the numerical target nor the admissible mass class can be fully specified from the supplied evidence alone.
\subsection{Short English statement}
The Ramos conjecture says that the least dimension allowing every collection of \(j\) admissible masses to be divided equally among all \(2^k\) sign regions of \(k\) affine hyperplanes equals the stated known lower bound, for all positive integers \(j,k\). The supplied record omits both the formula for that bound and the precise meaning of an admissible mass; the displayed measure formulation is therefore conditional.
\subsection{Sources}
\begin{itemize}
\item Supplied local UnsolvedMath registry record, permanent ID \(30002773\), exact source problem code \texttt{catalogue entry}, titled ``Ramos Conjecture for Equipartitions by Hyperplanes,'' in the Oberwolfach Reports Open Problems collection. The record names the Ramos conjecture but supplies no proposer attribution.
\item Source report: \emph{Geometric and Algebraic Combinatorics} (2015), Oberwolfach Reports, pp.~285--368, DOI \texttt{10.4171/owr/2015/5}. \url{https://doi.org/10.4171/owr/2015/5}. The supplied background describes Oberwolfach Reports as documenting mathematical workshops and their open-problem discussions.
\item Supplied literature reference: Haase--Blagojevic--Frick--Ziegler, \emph{Documenta Mathematica}. \url{https://doi.org/10.4171/DM/544}.
\item Supplied literature reference labeled ``primary arXiv version'': \url{https://arxiv.org/abs/1509.02959}.
\item Literature assessment dated exactly \texttt{2026-08-21}: the supplied curation labels the conjecture open in general and reports rigorous results for many parameter families and bounds, without equality with the lower bound for all \(j,k\). This is supplied dated curation derived from the OWR record, not independently verified here; its time-sensitive status and cited claims require independent verification.
\item The registry labels the statement \texttt{corrected\_verified}, with the explicit qualification that verification consisted only of safe notation, typography, and TeX-delimiter cleanup, preserving hypotheses, quantifiers, and subquestions. This label does not establish mathematical correctness.
\end{itemize}
\end{document}
